

% \section{Pola Perilaku Penyewa Kamar yang Berpotensi Melakukan Tindakan Terlarang}

% Tindakan terlarang dalam konteks penyewaan kamar mencakup berbagai aktivitas kriminal, khususnya penyalahgunaan narkoba dan praktik prostitusi. Penyalahgunaan narkoba mencakup penggunaan, penyimpanan, dan distribusi narkotika secara ilegal di dalam kamar sewa, termasuk berbagai jenis narkotika seperti metamfetamin, kokain, heroin, dan ekstasi sebagaimana diatur dalam Undang-Undang Nomor 35 Tahun 2009 tentang Narkotika. Sementara itu, prostitusi melibatkan eksploitasi seksual komersial yang dilakukan di kamar sewa, baik secara individu maupun dalam jaringan yang mana sering kali berkaitan dengan perdagangan manusia dan eksploitasi seksual anak (Pasal 296 dan 506 KUHP serta Undang-Undang Nomor 21 Tahun 2007 tentang Pemberantasan Tindak Pidana Perdagangan Orang).

% Regulasi hukum yang berlaku di Indonesia mencakup beberapa undang-undang utama yang berfungsi untuk mencegah aktivitas ilegal ini. Undang-Undang Nomor 35 Tahun 2009 tentang Narkotika mengatur larangan peredaran dan penggunaan narkotika secara ilegal, termasuk hukuman berat bagi pelaku. Pasal 296 dan 506 KUHP menetapkan sanksi bagi pihak yang memfasilitasi atau mengambil keuntungan dari prostitusi, sementara Undang-Undang Nomor 21 Tahun 2007 secara khusus melarang segala bentuk perdagangan manusia. Selain itu, beberapa daerah memiliki peraturan daerah (Perda) yang mengatur praktik penyewaan kamar guna mencegah tindak kriminal. 

% Di Kota Bandung, terdapat beberapa Peraturan Daerah (Perda) yang mengatur praktik penyewaan kamar dan bertujuan mencegah tindak kriminal. Salah satunya adalah Peraturan Daerah Kota Bandung Nomor 4 Tahun 2003 tentang Sewa Menyewa/Kontrak Rumah dan/atau Bangunan, yang mengatur ketentuan sewa menyewa atau kontrak rumah dan bangunan di Kota Bandung, termasuk kewajiban para pihak dan ketentuan lainnya yang bertujuan untuk menciptakan ketertiban dalam praktik penyewaan properti. Selain itu, terdapat Peraturan Daerah Kotamadya Daerah Tingkat II Bandung Nomor 07 Tahun 1986 tentang Izin Usaha Kepariwisataan yang mengatur tentang izin usaha kepariwisataan, termasuk losmen, penginapan remaja, pondok wisata, dan usaha rekreasi lainnya. Peraturan ini bertujuan untuk memastikan bahwa usaha-usaha tersebut beroperasi sesuai dengan ketentuan yang berlaku dan tidak disalahgunakan untuk aktivitas ilegal. Kemudian ada juga Peraturan Daerah Kota Bandung Nomor 9 Tahun 2019 tentang Ketertiban Umum, Ketentraman, dan Perlindungan Masyarakat yang mengatur berbagai aspek ketertiban umum di Kota Bandung, termasuk ketentuan yang dapat diterapkan untuk mencegah penyalahgunaan tempat tinggal atau kamar sewa untuk kegiatan kriminal. Dalam beberapa pasal, peraturan ini menetapkan larangan terhadap kegiatan yang mengganggu ketertiban umum yang dapat mencakup aktivitas ilegal di tempat sewa. Peraturan Daerah Kota Bandung Nomor 11 Tahun 2010 tentang Pelarangan, Pengawasan, dan Pengendalian Minuman Beralkohol mengatur larangan dan pengendalian peredaran minuman beralkohol di Kota Bandung. Meskipun fokus utamanya pada minuman beralkohol, penerapan peraturan ini dapat berkontribusi dalam mencegah tindak kriminal yang sering dikaitkan dengan konsumsi alkohol di tempat-tempat sewa. Terakhir, ada Peraturan Daerah Kota Bandung Nomor 13 Tahun 2019 tentang Perumahan, Kawasan Permukiman, dan Penanganan Kawasan Kumuh yang mengatur penyelenggaraan perumahan dan kawasan permukiman di Kota Bandung, termasuk upaya pencegahan dan peningkatan kualitas perumahan kumuh. Meskipun tidak secara spesifik mengatur penyewaan kamar, peraturan ini berkontribusi dalam menciptakan lingkungan permukiman yang tertib dan aman. Penerapan dan pengawasan terhadap Perda-Perda tersebut dilakukan oleh pemerintah daerah bersama aparat penegak hukum untuk memastikan bahwa praktik penyewaan kamar tidak disalahgunakan untuk kegiatan kriminal, termasuk di daerah Cihampelas dan wilayah lain di {Kota Bandung}.

% Penyewa yang berpotensi melakukan tindakan terlarang sering menunjukkan pola perilaku tertentu yang mencurigakan. Beberapa ciri utama termasuk penyewaan dalam jangka pendek dan tidak konsisten, sering berganti kamar atau lokasi dalam waktu yang tidak wajar, serta memiliki banyak tamu yang datang dan pergi dalam waktu singkat, terutama pada malam hari. Selain itu, pembayaran sering kali dilakukan dalam bentuk tunai dan tanpa identitas yang jelas, serta ada kecenderungan untuk menghindari interaksi dengan staf atau penghuni lain. Penyewa yang mencurigakan juga sering membawa banyak koper atau paket dalam kondisi mencurigakan, serta ditemukan alat-alat seperti jarum suntik, bong, atau bukti komunikasi digital terkait prostitusi (Leung, Yang, \& Dubin, 2018).

% Beberapa faktor yang menyebabkan hotel atau apartemen dipilih sebagai tempat tindakan terlarang antara lain kurangnya pengawasan dan keamanan, tidak adanya sistem identifikasi tamu yang ketat, serta lokasi yang strategis dengan akses mudah. Apartemen yang menerapkan sistem penyewaan fleksibel juga sering kali menjadi target utama karena tidak adanya batasan atau verifikasi bagi tamu yang sering berpindah-pindah kamar (Paraskevas \& Brookes, 2018). Selain itu, kurangnya kesadaran dan pelaporan dari pihak pengelola properti sering kali membuat aktivitas mencurigakan tidak terdeteksi.


% Untuk menangani penyewa yang melakukan tindakan kriminal, pengelola apartemen dan hotel dapat menerapkan berbagai strategi, termasuk penerapan sistem keamanan yang ketat dengan mewajibkan penyewa menunjukkan identitas asli yang diverifikasi serta menggunakan sistem CCTV dan patroli keamanan secara berkala. Pelatihan bagi staf dan manajemen sangat penting untuk mengenali pola perilaku mencurigakan serta menjalin kerja sama dengan pihak kepolisian dan organisasi terkait untuk edukasi (Zhang \& Potter, 2022). Selain itu, kerja sama dengan aparat penegak hukum perlu diperkuat dengan membentuk mekanisme pelaporan cepat dan meningkatkan koordinasi antara manajemen properti dengan kepolisian setempat. Teknologi juga dapat dimanfaatkan untuk mendeteksi aktivitas mencurigakan, misalnya melalui sistem analisis data dan {\it machine learning} yang mampu mengidentifikasi pola perilaku penyewa yang mencurigakan. Terakhir, penyuluhan dan kampanye kesadaran kepada penghuni apartemen dan staf mengenai bahaya aktivitas ilegal dapat membantu meningkatkan kewaspadaan dan mendorong pelaporan terhadap aktivitas mencurigakan.

% Hotel dan apartemen sering menjadi tempat yang dipilih untuk aktivitas ilegal seperti penyalahgunaan narkoba dan prostitusi karena berbagai faktor, termasuk lemahnya pengawasan dan regulasi penyewaan jangka pendek. Untuk mencegah dan menangani hal ini, diperlukan kombinasi strategi yang mencakup sistem keamanan yang ketat, pelatihan staf, kerja sama dengan aparat penegak hukum, serta pemanfaatan teknologi. Dengan pendekatan yang sistematis, pengelola properti dapat mengurangi risiko penyalahgunaan kamar sewa untuk aktivitas kriminal.








\section{Pola Perilaku Penyewa Kamar yang Berpotensi Melakukan Tindakan Terlarang}

Praktik penyewaan properti--seperti apartemen dan hotel-- rentan disalahgunakan untuk aktivitas ilegal, khususnya penyalahgunaan narkotika dan praktik prostitusi. Berdasarkan regulasi nasional, tindak pidana ini diatur secara ketat dan diancam dengan hukuman berat. Penyalahgunaan narkotika diatur dalam Undang-Undang Nomor 35 Tahun 2009 tentang Narkotika. Sementara itu, aktivitas yang memfasilitasi eksploitasi seksual komersial atau prostitusi dilarang keras berdasarkan Pasal 296 dan 506 KUHP, serta Undang-Undang Nomor 21 Tahun 2007 tentang Pemberantasan Tindak Pidana Perdagangan Orang (TPPO).

Di tingkat daerah, khususnya di Kota Bandung, pemerintah telah menetapkan berbagai instrumen hukum untuk memitigasi penyalahgunaan fungsi hunian dan menciptakan ketertiban umum. Beberapa regulasi utama tersebut sebagai berikut.
\begin{itemize}
    \item \textbf{Perda Kota Bandung No. 4 Tahun 2003} tentang Sewa Menyewa/Kontrak Rumah dan/atau Bangunan, yang mengatur kewajiban para pihak untuk menjaga ketertiban praktik penyewaan properti.
    \item \textbf{Perda Kotamadya Dati II Bandung No. 07 Tahun 1986} tentang Izin Usaha Kepariwisataan, guna memastikan penginapan dan losmen beroperasi sesuai izin dan tidak disalahgunakan.
    \item \textbf{Perda Kota Bandung No. 11 Tahun 2010} tentang Pelarangan, Pengawasan, dan Pengendalian Minuman Beralkohol, yang membatasi pemicu tindak kriminal di tempat sewa.
    \item \textbf{Perda Kota Bandung No. 9 Tahun 2019} (Ketertiban Umum) dan \textbf{Perda No. 13 Tahun 2019} (Perumahan dan Kawasan Permukiman), yang memberikan dasar hukum penindakan terhadap aktivitas yang mengganggu ketentraman masyarakat di area permukiman.
\end{itemize}

Menurut literatur kriminologi dan perhotelan, penyewa yang berpotensi melakukan tindakan terlarang umumnya menunjukkan pola perilaku (\textit{behavioral patterns}) spasial dan temporal yang anomali. Leung, Yang, dan Dubin (2018) mengidentifikasi beberapa indikator perilaku mencurigakan, antara lain penyewaan dalam jangka waktu sangat pendek (\textit{micro-stays}) yang sering dilakukan berulang; intensitas keluar-masuk tamu yang tinggi dalam waktu singkat, terutama pada malam hari; kecenderungan membayar menggunakan uang tunai untuk menghindari pelacakan identitas; serta keengganan berinteraksi dengan staf atau penghuni lain (\textit{avoidance behavior}). Selain itu, kecurigaan fisik sering kali didukung dengan penemuan barang bawaan tak lazim atau residu aktivitas ilegal (seperti jarum suntik atau bong).

Tingginya kerentanan apartemen terhadap aktivitas ini disebabkan oleh karakteristik struktural penyewaan itu sendiri. Paraskevas dan Brookes (2018) dalam kajian keamanan perhotelan menyatakan bahwa properti yang menerapkan sistem sewa fleksibel menjadi target utama karena lemahnya verifikasi identitas, minimnya pengawasan langsung (\textit{blind spots}), serta tingkat privasi yang tinggi yang menyulitkan deteksi dini oleh pengelola. 

Untuk menangani ancaman tersebut, pengelola properti memerlukan pendekatan keamanan yang berlapis (\textit{multilayered security}). Pendekatan ini mencakup penerapan verifikasi identitas otentik, optimalisasi pengawasan CCTV, dan pelatihan staf untuk mengenali anomali perilaku tamu (Zhang \& Potter, 2022). Lebih lanjut, adopsi teknologi berbasis data dan \textit{machine learning} menjadi esensial di era modern. Sistem deteksi otomatis dapat menganalisis rekam jejak transaksi, durasi inap, dan pola kedatangan penyewa untuk menghasilkan peringatan dini (\textit{early warning system}) bagi pihak manajemen. Melalui kombinasi antara protokol keamanan fisik, kolaborasi dengan aparat penegak hukum, dan sistem algoritma cerdas, risiko penyalahgunaan kamar sewa untuk aktivitas kriminal dapat direduksi secara signifikan.